\documentclass[11pt]{article}
\usepackage[margin=25mm]{geometry}
\usepackage{fontspec}
\setmainfont{Latin Modern Roman}
\usepackage{amsmath,amssymb,mathtools}
\usepackage{unicode-math}
\setmathfont{Latin Modern Math}
\usepackage{xurl,hyperref}
\hypersetup{colorlinks=true,urlcolor=blue}
\setcounter{secnumdepth}{0}
\setlength{\parindent}{0pt}
\setlength{\parskip}{5pt}
\setlength{\emergencystretch}{3em}
\newcommand{\sref}[2]{\hyperlink{source:#1:#2}{[S#2]}}
\newcommand{\cataloguescope}[1]{\paragraph{Recorded target.}#1}
\newcommand{\R}{\mathbb R}
\newcommand{\C}{\mathbb C}
\newcommand{\N}{\mathbb N}
\newcommand{\Z}{\mathbb Z}
\newcommand{\Q}{\mathbb Q}
\newcommand{\F}{\mathbb F}
\newcommand{\eps}{\varepsilon}
\begin{document}
% UnsolvedMath ID 30001819; source catalogue code OWR-5160-001
\setcounter{section}{30001818}
\section{Energy Cost of Density Defects in Fermi Seas}\label{30001819}
{\small UnsolvedMath ID 30001819 · Mathematical Physics}
\subsection{Definitions and mathematical statement}

\paragraph{Shared notation.}$\mathbb N=\{1,2,\ldots\}$; the standard number fields have their usual meanings. All additional conventions are stated locally.


\paragraph{Definitions and conventions.} For a free Fermi sea in $\mathbb R^d$, chemical potential $\mu>0$ and spin degeneracy $q$ fix $\rho_0=q(2\pi)^{-d}|\{k:|k|^2\le\mu\}|$. Its one-particle density matrix is $\Pi_-=\mathbf1_{(-\infty,\mu)}(-\Delta)$. An admissible fermionic state has $0\le\gamma\le1$, density $\rho_\gamma$, and relative perturbation $Q=\gamma-\Pi_-$. The density defect $\delta\rho=\rho_\gamma-\rho_0$ can be negative but satisfies $\delta\rho\ge-\rho_0$. Relative kinetic energy is the properly defined relative trace of $(-\Delta-\mu)Q$, not the difference of two divergent absolute energies. A positive-density Lieb--Thirring bound compares that energy to an integral of the convex remainder $\rho_\gamma^{1+2/d}-\rho_0^{1+2/d}-(1+2/d)\rho_0^{2/d}\delta\rho$.

\paragraph{Source scope and qualifications.} The official report already proves a positive-density lower bound in d>=2. The remaining optimal-cost or sharp-constant task is more specific than the catalogue's title-only dated assessment. Dimensions, spin normalization, finite-relative-energy domain, and density admissibility must be retained; a trace of two infinite energies is not a valid definition.

\paragraph{Catalogue mathematical target.}
\paragraph{Catalogue wording retained for traceability.}
Determine the optimal energy cost, or sharp Lieb–Thirring-type lower bound, for creating a density perturbation including a hole in an infinitely extended free Fermi sea of prescribed positive density $\rho_0$.

\paragraph{Saved literature context.} The supplied catalogue's assessment dated 2026-08-21 records: This record is only a title and contains no mathematically testable assertion or resolution; it is therefore retained as open rather than inferred solved. The assessment is retained as dated provenance; it is not a new comprehensive literature or open-status certification.
\subsection{Short English statement}
Find the least energy needed to make particle-density excesses or holes in an infinite Fermi sea, including sharp lower-bound constants.
\subsection{Research attempt}
\paragraph{Model, date and investigated scope.}
OpenAI GPT-6 Astra (configured model \texttt{gpt-6-astra}), ultra reasoning; generated 11 October 2026. Computes the semiclassical homogeneous convex remainder, small-defect coefficient and complete-hole cost with spin normalization.

\paragraph{Sources and literature check.}
The complete supplied definition, including source qualifications and the dated literature assessment, was reviewed. The calculation below is an elementary restricted attempt motivated by [S2]; it does not claim novelty or a fresh literature-status determination.

\paragraph{Context and investigated scope.} The complete supplied definition, including source qualifications and the dated literature assessment, was reviewed. The calculation below is an elementary restricted attempt motivated by [S2]; it does not claim novelty or a fresh literature-status determination.

\paragraph{The homogeneous density-cost benchmark.} Let $v_d$ be the volume of the unit ball and $q$ the spin degeneracy. Filling a momentum ball of radius $k_F$ gives $\rho=qv_dk_F^d/(2\pi)^d$. Direct radial integration of $|k|^2$ gives kinetic energy density
\[
 e(\rho)=K_{d,q}\rho^{1+2/d},\qquad
 K_{d,q}=\frac d{d+2}\left(\frac{(2\pi)^d}{qv_d}\right)^{2/d}.
\]
The derivative at $\rho_0$ is the Fermi chemical potential $\mu=k_{F,0}^2$. Hence the homogeneous grand-canonical cost of changing density by $s\geq-\rho_0$ is
\[
 e(\rho_0+s)-e(\rho_0)-\mu s
 =K_{d,q}\bigl((\rho_0+s)^p-\rho_0^p-p\rho_0^{p-1}s\bigr),\quad p=1+2/d.
\]
Strict convexity makes it nonnegative, with quadratic small-$s$ coefficient $K_{d,q}p(p-1)\rho_0^{p-2}/2$.

\paragraph{The complete-hole boundary.} At $s=-\rho_0$, the expression is $K_{d,q}(p-1)\rho_0^p>0$. Thus removing all local particles need not lower the relative energy: subtraction of the chemical-potential term matters. At $s=0$ it vanishes to second order as required.

\paragraph{Sharpness gap.} These are uniform momentum fillings, whose infinite-volume total relative energy is not finite; they provide a cost per volume benchmark only. Turning them into admissible localized defects requires boundary control, and proving a sharp lower bound for arbitrary $0\leq\gamma\leq1$ requires quantum spatial estimates. The source already proves a nonsharp positive-density bound in $d\geq2$. This elementary calculation neither replaces it nor establishes the optimal constant or optimal density-hole profile.

\paragraph{Verification and adversarial audit.}
Checked phase-space normalization, derivative equal to chemical potential, hole endpoint and distinction between energy density and a finite relative trace. Model self-audit only; no independent human proof checking. Prompts A and B were applied to this canonical record and its qualified source scope. Checked phase-space normalization, derivative equal to chemical potential, hole endpoint and distinction between energy density and a finite relative trace. Retained unresolved status and zero original-target progress; Localized finite-relative-energy constructions and a sharp universal quantum lower bound are missing.

Coordinator checked the momentum-ball kinetic constant, chemical potential and complete-hole relative energy, with localization and optimality left open. This additional check was model-assisted, not independent human verification.

\paragraph{Tools, process and contributions.}
Python 3 was used for file preparation and structural checks; no numerical experiment or theorem prover was used. Contribution prompt A was followed by prompt B. GPT-6 Astra with ultra reasoning drafted the research and performed its stated self-audit; any additional coordinator contribution is identified in the verification and audit. A GPT-6-based Codex coordinator assembled the manuscript and checked repository integration; its exact variant and reasoning configuration were not exposed. No human mathematical contribution, independent human proof check or formal proof certification is supplied. Complete derivations and the consulted source references are included here; private reasoning transcripts are not published.

\paragraph{Final status and first remaining gap.}
\textbf{UNRESOLVED}. Localized finite-relative-energy constructions and a sharp universal quantum lower bound are missing.

\section{Completion Estimate}
\noindent\textbf{COMPLETION: 0\%}
The calculation is an auxiliary or restricted attempt and establishes no newly verified portion of the full source target. Localized finite-relative-energy constructions and a sharp universal quantum lower bound are missing.

\subsection{Sources}
\begin{itemize}
\item[\hypertarget{source:30001819:1}{S1}] ULAM.AI and the UnsolvedMath Contributors, UnsolvedMath, record 30001819 (OWR-5160-001). Catalogue curation under CC BY 4.0; definitions and English exposition added locally.
\url{https://huggingface.co/datasets/ulamai/UnsolvedMath}.
\item[\hypertarget{source:30001819:2}{S2}] Mathematical Methods in Quantum Chemistry (2011), Oberwolfach Reports, DOI 10.4171/owr/2011/32, pp. 1769--1843; official report source and mathematical context.
\url{https://doi.org/10.4171/owr/2011/32}.
\item[\hypertarget{source:30001819:3}{S3}] Official EMS Press PDF of the same Oberwolfach report; PDF pages 23--24, printed pages 1791--1792.
\url{https://ems.press/content/serial-article-files/46348}.
\end{itemize}
\label{end:30001819}
\end{document}
